\section{Residue characteristic polynomials and short algebraic certificates}
\label{sec:residue}

The Potts improvements act on a scalar obstruction. A separate line of
progress concerns the local rings that occur in exact categorical
Khovanov computation. Earlier ProveIt work already uses nilpotent radical
perturbations and finite geometric-series inverses
\cite{report08,report12}. The immediately preceding continuation constructs
Fitting decompositions of supplied \emph{scalar} binary endomorphisms.
We now extend the algebraic certificate mechanism to supplied ring-valued,
same-matching blocks.

The ingredients are standard Cayley--Hamilton, scalar Frobenius, polynomial
Chinese remainders, and Fitting decomposition. The theorem is stated with a
complete elementary proof because the applicable ring and matrix hypotheses
matter. This is not a literature-wide novelty claim about matrix nil indices.
No production scanner speedup from this supplemental result is claimed.

\subsection{A residue-only annihilator}

Fix a prime $p$ and positive integers $b,r$. Set
\[
 R_{b,p}=\F_p[x_1,\ldots,x_b]/(x_1^p,\ldots,x_b^p),\qquad
 J=(x_1,\ldots,x_b),\qquad B=b(p-1).
\]
Let $\rho:R_{b,p}\to\F_p$ take the constant coefficient.
The monomial basis has exponents $0\le\alpha_i<p$, so
$J^{B+1}=0$ and $J^B\ne0$. Scalar Frobenius gives
\begin{equation}
 a^p=\rho(a)\qquad(a\in R_{b,p}).
 \label{eq:scalar-frobenius}
\end{equation}
This identity concerns scalar elements of a commutative ring. The map
$A\mapsto A^p$ on matrices is not entrywise Frobenius and need not be additive.

\begin{theorem}[Residue characteristic-polynomial certificate]
\label{thm:residue}
For $A\in M_r(R_{b,p})$, put $A_0=\rho(A)$ entrywise and
$\chi_0(t)=\det(tI-A_0)\in\F_p[t]$. Then
\[
 \boxed{\chi_0(A)^p=0.}
\]
Thus $A$ has an explicitly computable annihilating polynomial
$\chi_0(t)^p$ of degree $pr$, independent of $b$.
\end{theorem}
\begin{proof}
Write $\chi_A(t)=\sum_{i=0}^r c_it^i$ over $R_{b,p}$, with $c_r=1$.
Taking constant coefficients commutes with the determinant, so
$\chi_0(t)=\sum_i\rho(c_i)t^i$.
Cayley--Hamilton over a commutative ring gives $\sum_i c_iA^i=0$.
These summands commute with each other: each is a scalar multiple of a
power of the same matrix. Consequently raising their sum to its $p$th power
is legitimate and yields
\[
 0=\sum_i c_i^p A^{pi}
   =\sum_i\rho(c_i)A^{pi}
   =(\chi_0(t)^p)(A)=\chi_0(A)^p.
\]
No Frobenius identity for arbitrary noncommuting matrix summands was used.
\end{proof}

The practical case here is $p=2$, where the dot algebra of a same-matching
block has square-zero variables. Its annihilator degree is at most $2r$,
even if the scalar coefficient space has dimension $2^b$.

\subsection{A sharp radical-matrix bound}

\begin{theorem}[Sharp universal radical cutoff]\label{thm:radical}
Every $N\in M_r(J)$ satisfies
\[
 N^L=0,\qquad L=\min\{b(p-1)+1,pr\}.
\]
For every $b,r\ge1$ and prime $p$, there is such an $N$ with
$N^{L-1}\ne0$.
\end{theorem}
\begin{proof}
The ideal bound gives $N^{B+1}=0$. Since $N_0=0$,
$\chi_0(t)=t^r$, and Theorem~\ref{thm:residue} gives $N^{pr}=0$.

For sharpness, let $s=\min(b,r)$. Use an $s\times s$ upper-bidiagonal block
with distinct variables $x_1,\ldots,x_s$ on the diagonal. Of its $s-1$
superdiagonal positions, label
\[
 e=\min\{s-1,(b-s)(p-1)\}
\]
with variables among $x_{s+1},\ldots,x_b$, each used at most $p-1$ times.
Label the remaining positions with diagonal variables, again with each
used at most $p-1$ times. There is enough capacity since
$s-1\le s(p-1)$. Let $a_i$ be the number of superdiagonal entries labeled
$x_i$ for $i\le s$.

In a power of this bidiagonal matrix, a path from row $1$ to column $s$
traverses every superdiagonal position in its forced order and can wait
at vertex $i$ by using the diagonal $x_i$. Choose exactly $p-1-a_i$ waits
at that vertex. The resulting monomial has degree
\[
 D=(s-1)+\sum_{i=1}^s(p-1-a_i)
  =s(p-1)+e
  =\min\{b(p-1),pr-1\}=L-1.
\]
Every variable exponent is below $p$. Its coefficient is one: the
superdiagonal order is forced and the distinct diagonal variables specify
the waiting counts uniquely. Hence it survives in the $(1,s)$ entry of
$N^D$. Embed the block into an $r\times r$ zero matrix when $s<r$.
\end{proof}

For $p=2$ the cutoff is $\min(b+1,2r)$.
When $b\ge2r-1$, a particularly transparent witness has $r$ distinct
diagonal variables and $r-1$ further distinct superdiagonal variables.
The $(1,r)$ entry of $N^{2r-1}$ is their product.
This shows that replacing $2r$ by $r$, based on a scalar square-zero
intuition, would be wrong.

If instead $A_0$ is nilpotent of index $s$, then $A^s\in M_r(J)$.
Combining the ideal bound with Theorem~\ref{thm:residue} yields
\begin{equation}
 A^{\min\{pr,\ s[b(p-1)+1]\}}=0.
 \label{eq:nilpotent-residue}
\end{equation}
The $pr$ term can be attained already with one variable: the companion
matrix of $t^r-x_1$ has $A^r=x_1I$ and nilpotence index $pr$.
We do not claim simultaneous sharpness of \eqref{eq:nilpotent-residue}
for every prescribed triple $(b,r,s)$.

\subsection{Inverses without a determinant over the large ring}

\begin{corollary}\label{cor:inverse}
For $N\in M_r(J)$,
\[
 (I+N)^{-1}=\sum_{j=0}^{L-1}(-N)^j,\qquad
 L=\min\{b(p-1)+1,pr\}.
\]
More generally, if $A_0$ is invertible and
$\chi_0(t)=\sum_{i=0}^r c_it^i$, then
\[
 A^{-1}=-c_0^{-1}\sum_{i=1}^r c_iA^{pi-1}.
\]
\end{corollary}
\begin{proof}
The first formula telescopes using $N^L=0$.
For the second, Theorem~\ref{thm:residue} says
$c_0I+\sum_{i=1}^r c_iA^{pi}=0$.
Multiplying the proposed inverse on either side by $A$ gives $I$.
\end{proof}

At $p=2$, the minus signs disappear and $c_0=1$.
The characteristic polynomial is computed over the small base field, not
by expanding a determinant with $2^b$ coefficients in each entry.
The powers of $A$ still require actual ring arithmetic.
If a geometric series is evaluated by doubling a power-and-partial-sum
pair, it takes $O(\log L)$ matrix products. The comparison with earlier
methods must acknowledge this: reducing a cutoff from $b+1$ to $2r$
does not automatically save $b/r$ matrix products in an already doubled
implementation.

\subsection{Constructive Fitting and Drazin data}

Factor the residue characteristic polynomial as
\[
 \chi_0(t)=t^a q_0(t),\qquad q_0(0)\ne0.
\]
The notation $q_0$ here is a polynomial and is unrelated to a Potts color
count. Polynomial CRT gives a residue class $e(t)$ modulo
$t^{pa}q_0(t)^p$ satisfying
\[
 e(t)\equiv0\pmod{t^{pa}},\qquad
 e(t)\equiv1\pmod{q_0(t)^p}.
\]
Choose its representative of degree below $pr$ and set $E=e(A)$.

\begin{theorem}[Explicit Fitting projectors]\label{thm:fitting}
The operators above satisfy
\[
 E^2=E,\quad AE=EA,\quad A^{pa}(I-E)=0,\quad q_0(A)^pE=0.
\]
There is a polynomially specified $D$ with
$AD=DA=E$ and $DE=ED=D$.
The modules $\ker E$ and $\im E$ are free of ranks $a$ and $r-a$,
and a concrete basis change can be lifted from residue projector columns.
\end{theorem}
\begin{proof}
All displayed identities for $E$ follow from polynomial divisibility by
$t^{pa}q_0^p$ and Theorem~\ref{thm:residue}.
Choose $z(t)$ by CRT with $z=0$ modulo $t^{pa}$ and
$z=t^{-1}$ modulo $q_0^p$, and let $D=z(A)$.
The claimed identities follow by the same argument.
In particular $A$ is nilpotent on $\ker E$ and invertible on $\im E$.
The limiting cases $a=0$ and $a=r$ use $E=I$ and $E=0$ respectively.

Choose columns of $E_0$ spanning $\im E_0$ over $\F_p$, and columns of
$I-E_0$ spanning $\ker E_0$. Take the corresponding columns of $E$ and
$I-E$ and arrange the kernel columns first in a matrix $P$.
Its residue $P_0$ is invertible. Since $R_{b,p}$ is local with maximal
ideal $J$, $\det P$ is a unit, so $P$ is invertible.
The column identities give
$EP=P\operatorname{diag}(0_a,I_{r-a})$, proving the rank and basis claims.
\end{proof}

The resulting $D$ is the Fitting/Drazin inverse:
$DAD=D$ and $A^{pa+1}D=A^{pa}$.
The method constructs the decomposition inside matrices over $R_{b,p}$.
It need not first expand the module to dimension $rp^b$ over the base field.
It nevertheless must represent and manipulate its ring coefficients.

\subsection{Intertwining across categorical blocks}

Suppose a finite additive $\F_p$-linear category contains same-matching
objects $X_i$ with $\End(X_i)=R_{b_i,p}$, and a degree-preserving chain
endomorphism $Q$ is block diagonal by degree and matching.
Write $A_i$ for its block on $X_i^{r_i}$. A differential component $T$
satisfies $A_jT=TA_i$ under actual categorical composition.

Construct the respective $E_i,D_i$ and $E_j,D_j$ by the preceding theorem.
Their projector polynomials need not be the same.
For $N$ large enough to kill both nilpotent restrictions,
\[
 (I-E_j)TE_i
  =(I-E_j)TA_i^ND_i^N
  =(I-E_j)A_j^NTD_i^N=0.
\]
Similarly,
\[
 E_jT(I-E_i)
  =D_j^NA_j^NT(I-E_i)
  =D_j^NTA_i^N(I-E_i)=0.
\]
Hence $E_jT=TE_i$, and the nilpotent and invertible images form
complementary subcomplexes. The lifted bases exhibit a genuine chain
isomorphism when every inverse and categorical cross term is evaluated.

The characteristic-polynomial theorem is used only inside the commutative
same-matching endomorphism blocks. Applying it directly to an arbitrary
mixed-matching matrix with noncommutative cobordism entries would be invalid.
Finding useful $Q$, keeping it small, and exploiting the resulting summands
are separate algorithmic problems.

\subsection{How this advances the previous continuation}

The previous scalar Fitting procedure already stabilizes in at most its
matrix dimension over $\F_2$. The $2r$ annihilator does not accelerate that
scalar computation. Its contribution is a ring-valued certificate whose
polynomial degree is controlled by block multiplicity rather than only by
the number of dot variables. The sharpness theorem prevents an unjustified
further truncation.

The supplied oracle represents monomials by exponent tuples and checks the
identities independently of the production dot algebra. Its characteristic
polynomial routines are intended for small exhaustive experiments, not as a
scalable implementation. The exact counts are reported in
Section~\ref{sec:experiments}. No bound on all scanner multiplicities,
all block sizes, or closure of a succinct representation follows from
these local theorems.
